\section{Tutorial 8: Training Ablations (BatchNorm, Dropout, Residuals)}

\subsection{Objectives}

This tutorial studies how architectural and training choices affect learning:

\begin{itemize}
    \item Understand the role of normalization, regularization, and residual connections
    \item Compare training dynamics under different configurations
    \item Develop intuition through controlled ablation experiments
\end{itemize}

\medskip

\noindent
\textbf{Focus.}  
Identify which components make deep networks trainable and generalize well.

\subsection{Exercise 1: What is an Ablation Study?}

\textbf{Question.}  
What is the purpose of an ablation study in machine learning?

\medskip

\noindent
\textbf{Discussion.}

\begin{itemize}
    \item Systematically remove or modify components
    \item Measure their effect on performance
\end{itemize}

\medskip

\noindent
\textbf{Insight.}  
Ablation reveals causal contributions of design choices.

\subsection{Exercise 2: Baseline Model}

Consider a deep fully connected network:

\begin{itemize}
    \item depth: $L = 10$
    \item activation: ReLU
    \item optimizer: SGD
\end{itemize}

\medskip

\noindent
\textbf{Task.}  
Predict training behavior without:
\begin{itemize}
    \item normalization
    \item residual connections
\end{itemize}

\medskip

\noindent
\textbf{Expected outcome.}

\begin{itemize}
    \item unstable gradients
    \item slow or failed convergence
\end{itemize}

\medskip

\noindent
\textbf{Insight.}  
Depth alone does not guarantee learnability.

\subsection{Exercise 3: Effect of Batch Normalization}

Modify the baseline:

\begin{itemize}
    \item add batch normalization after each layer
\end{itemize}

\medskip

\noindent
\textbf{Tasks.}
\begin{enumerate}
    \item Predict changes in training speed
    \item Predict effect on gradient stability
\end{enumerate}

\medskip

\noindent
\textbf{Expected outcome.}

\begin{itemize}
    \item faster convergence
    \item more stable gradients
\end{itemize}

\medskip

\noindent
\textbf{Insight.}  
Normalization stabilizes optimization.

\subsection{Exercise 4: Effect of Dropout}

Modify the baseline:

\begin{itemize}
    \item add dropout with rate $p$
\end{itemize}

\medskip

\noindent
\textbf{Tasks.}
\begin{enumerate}
    \item Predict effect on training loss
    \item Predict effect on test performance
\end{enumerate}

\medskip

\noindent
\textbf{Expected outcome.}

\begin{itemize}
    \item slower training (higher training loss)
    \item improved generalization
\end{itemize}

\medskip

\noindent
\textbf{Insight.}  
Regularization trades training performance for robustness.

\subsection{Exercise 5: Effect of Residual Connections}

Modify the baseline:

\begin{itemize}
    \item replace layers with residual blocks:
    \[
    h_{k+1} = h_k + F(h_k)
    \]
\end{itemize}

\medskip

\noindent
\textbf{Tasks.}
\begin{enumerate}
    \item Predict training behavior
    \item Compare to plain deep network
\end{enumerate}

\medskip

\noindent
\textbf{Expected outcome.}

\begin{itemize}
    \item easier optimization
    \item deeper networks become trainable
\end{itemize}

\medskip

\noindent
\textbf{Insight.}  
Residual connections improve gradient flow structurally.

\subsection{Exercise 6: Combined Model}

Consider a model with:

\begin{itemize}
    \item ReLU activation
    \item batch normalization
    \item residual connections
    \item optional dropout
\end{itemize}

\medskip

\noindent
\textbf{Tasks.}
\begin{enumerate}
    \item Predict overall training behavior
    \item Explain why this combination is effective
\end{enumerate}

\medskip

\noindent
\textbf{Expected outcome.}

\begin{itemize}
    \item fast convergence
    \item stable gradients
    \item good generalization
\end{itemize}

\medskip

\noindent
\textbf{Insight.}  
Modern deep learning relies on combining multiple techniques.

\subsection{Exercise 7: Comparative Table}

Fill in the following table:

\[
\begin{array}{c|c|c|c}
\text{Model} & \text{Training Stability} & \text{Speed} & \text{Generalization} \\
\hline
\text{Baseline} & & & \\
\text{+ BN} & & & \\
\text{+ Dropout} & & & \\
\text{+ Residual} & & & \\
\text{Combined} & & & \\
\end{array}
\]

\medskip

\noindent
\textbf{Discussion.}

Students should observe:

\begin{itemize}
    \item BN $\rightarrow$ stability + speed
    \item Dropout $\rightarrow$ generalization
    \item Residual $\rightarrow$ depth scalability
\end{itemize}

\subsection{Exercise 8 (Discussion): Design Principles}

\textbf{Question.}  
What general principles can we extract from these experiments?

\medskip

\noindent
\textbf{Discussion points.}

\begin{itemize}
    \item Control signal propagation (BN, residuals)
    \item Prevent overfitting (dropout)
    \item Combine techniques for best performance
\end{itemize}

\medskip

\noindent
\textbf{Insight.}  
Successful deep learning is the result of carefully engineered systems.

\subsection{Summary}

\begin{itemize}
    \item Batch normalization stabilizes training
    \item Dropout improves generalization
    \item Residual connections enable depth
    \item Combined techniques yield robust models
\end{itemize}

\medskip

\noindent
\textbf{Preparation for next lecture.}  
We now study convolutional neural networks, which incorporate structural assumptions about data into architecture design.